% ============================================================================
%  PRÉAMBULE BEAMER — copié du projet « Cours université physique médicale »
%  (commun/preambule-presentation.tex), polices versionnées dans
%  commun/polices/, palette RX. Modifications locales : chemin du .bib,
%  macros.tex non chargé (il renomme MTF en FTM, l'exposé garde les sigles
%  anglais).
%  Le document a déjà déclaré \documentclass[aspectratio=169, 11pt]{beamer}.
% ============================================================================

\usepackage{etoolbox}
\usepackage{xcolor}

% --- Polices (fontspec, chemin relatif au dépôt) ---
\usepackage[no-math]{fontspec}
\defaultfontfeatures{Ligatures=TeX}
\setsansfont{Poppins}[Extension=.ttf,
  UprightFont=*-Regular, BoldFont=*-Bold, ItalicFont=*-Italic, BoldItalicFont=*-BoldItalic]
\setmainfont{Poppins}[Extension=.ttf,
  UprightFont=*-Regular, BoldFont=*-Bold, ItalicFont=*-Italic, BoldItalicFont=*-BoldItalic]
\newfontfamily\lightfont{Poppins-Light}[Extension=.ttf]
\newfontfamily\mediumfont{Poppins-Medium}[Extension=.ttf]

% --- Mathématiques : chiffres et lettres latines en Poppins (mathastext),
%     grec en Inter, symboles en Noto Sans Math ---
\usepackage{mathtools}
\usepackage[italic]{mathastext}
\newfontfamily\intermath{Inter}[Extension=.otf,
  UprightFont=*-Regular, BoldFont=*-Bold, ItalicFont=*-Italic, BoldItalicFont=*-BoldItalic]
\newfontfamily\notosansmath{NotoSansMath-Regular}[Extension=.ttf]
\makeatletter
{\intermath\xdef\InterFamily{\f@family}}
{\notosansmath\xdef\NotoMathFamily{\f@family}}
\makeatother
\DeclareSymbolFont{greeklower}{TU}{\InterFamily}{m}{it}
\SetSymbolFont{greeklower}{bold}{TU}{\InterFamily}{b}{it}
\DeclareSymbolFont{greekupper}{TU}{\InterFamily}{m}{n}
\SetSymbolFont{greekupper}{bold}{TU}{\InterFamily}{b}{n}
\DeclareSymbolFont{mathsans}{TU}{\NotoMathFamily}{m}{n}
% Grec minuscule (italique, classe 7)
\Umathchardef\alpha=7 \symgreeklower "03B1 \Umathchardef\beta=7 \symgreeklower "03B2
\Umathchardef\gamma=7 \symgreeklower "03B3 \Umathchardef\delta=7 \symgreeklower "03B4
\Umathchardef\epsilon=7 \symgreeklower "03F5 \Umathchardef\varepsilon=7 \symgreeklower "03B5
\Umathchardef\zeta=7 \symgreeklower "03B6 \Umathchardef\eta=7 \symgreeklower "03B7
\Umathchardef\theta=7 \symgreeklower "03B8 \Umathchardef\vartheta=7 \symgreeklower "03D1
\Umathchardef\iota=7 \symgreeklower "03B9 \Umathchardef\kappa=7 \symgreeklower "03BA
\Umathchardef\lambda=7 \symgreeklower "03BB \Umathchardef\mu=7 \symgreeklower "03BC
\Umathchardef\nu=7 \symgreeklower "03BD \Umathchardef\xi=7 \symgreeklower "03BE
\Umathchardef\pi=7 \symgreeklower "03C0 \Umathchardef\varpi=7 \symgreeklower "03D6
\Umathchardef\rho=7 \symgreeklower "03C1 \Umathchardef\varrho=7 \symgreeklower "03F1
\Umathchardef\sigma=7 \symgreeklower "03C3 \Umathchardef\varsigma=7 \symgreeklower "03C2
\Umathchardef\tau=7 \symgreeklower "03C4 \Umathchardef\upsilon=7 \symgreeklower "03C5
\Umathchardef\phi=7 \symgreeklower "03D5 \Umathchardef\varphi=7 \symgreeklower "03C6
\Umathchardef\chi=7 \symgreeklower "03C7 \Umathchardef\psi=7 \symgreeklower "03C8
\Umathchardef\omega=7 \symgreeklower "03C9
% Grec majuscule (droit, classe 0)
\Umathchardef\Gamma=0 \symgreekupper "0393 \Umathchardef\Delta=0 \symgreekupper "0394
\Umathchardef\Theta=0 \symgreekupper "0398 \Umathchardef\Lambda=0 \symgreekupper "039B
\Umathchardef\Xi=0 \symgreekupper "039E \Umathchardef\Pi=0 \symgreekupper "03A0
\Umathchardef\Sigma=0 \symgreekupper "03A3 \Umathchardef\Upsilon=0 \symgreekupper "03A5
\Umathchardef\Phi=0 \symgreekupper "03A6 \Umathchardef\Psi=0 \symgreekupper "03A8
\Umathchardef\Omega=0 \symgreekupper "03A9
% Relations (classe 3), opérateurs binaires (2), ordinaires (0)
\Umathchardef\approx=3 \symmathsans "2248 \Umathchardef\sim=3 \symmathsans "223C
\Umathchardef\simeq=3 \symmathsans "2243 \Umathchardef\propto=3 \symmathsans "221D
\Umathchardef\leq=3 \symmathsans "2264 \Umathchardef\geq=3 \symmathsans "2265
\Umathchardef\neq=3 \symmathsans "2260 \Umathchardef\equiv=3 \symmathsans "2261
\Umathchardef\ll=3 \symmathsans "226A \Umathchardef\gg=3 \symmathsans "226B
\Umathchardef\rightarrow=3 \symmathsans "2192 \Umathchardef\leftarrow=3 \symmathsans "2190
\Umathchardef\Rightarrow=3 \symmathsans "21D2 \Umathchardef\Leftarrow=3 \symmathsans "21D0
\Umathchardef\leftrightarrow=3 \symmathsans "2194
\Umathchardef\times=2 \symmathsans "00D7 \Umathchardef\pm=2 \symmathsans "00B1
\Umathchardef\mp=2 \symmathsans "2213 \Umathchardef\cdot=2 \symmathsans "22C5
\Umathchardef\circ=2 \symmathsans "2218
\Umathchardef\infty=0 \symmathsans "221E \Umathchardef\partial=0 \symmathsans "2202
\let\le\leq \let\ge\geq \let\ne\neq \let\to\rightarrow \let\gets\leftarrow

% --- Langue, unités, graphiques ---
\usepackage[french]{babel}
\usepackage[autostyle=true, french=guillemets]{csquotes}
\usepackage[locale=FR, per-mode=symbol, separate-uncertainty=true, range-phrase={ à }, range-units=single]{siunitx}
\usepackage{booktabs}
\usepackage{graphicx}
\usepackage{tikz, pgfplots}
\pgfplotsset{compat=1.18}
\usetikzlibrary{calc, positioning, arrows.meta, shapes.geometric, decorations.pathreplacing, decorations.pathmorphing, backgrounds}
\usepackage[most]{tcolorbox}
\usepackage{appendixnumberbeamer}
\usepackage[backend=biber, style=authoryear, maxcitenames=2, giveninits=true, uniquename=init]{biblatex}
\addbibresource{../references.bib}
\usepackage{hyperref}
\hypersetup{colorlinks=false, pdfborderstyle={/S/U/W 0}}

\input{\commun/palettes}
\input{\commun/boites}

% ============================================================================
%  RÉGLAGES BEAMER
% ============================================================================
\usetheme{default}
\usecolortheme{default}
\usefonttheme{professionalfonts}
\setbeamertemplate{navigation symbols}{}
\setbeamertemplate{headline}{}
\setbeamertemplate{frametitle continuation}{}
\setbeamertemplate{blocks}[default]
\setbeamertemplate{caption}{\raggedright\insertcaption\par}
\setbeamerfont{caption}{size=\scriptsize}
\setbeamercolor{caption}{fg=StoneLight}

\setbeamercolor{normal text}{fg=Stone, bg=SoftWhite}
\setbeamercolor{structure}{fg=Slate}
\setbeamercolor{alerted text}{fg=WarmAmber}
\setbeamercolor{example text}{fg=Sage}
\setbeamercolor{title}{fg=Slate}
\setbeamercolor{subtitle}{fg=Sage}
\setbeamercolor{author}{fg=Stone}
\setbeamercolor{institute}{fg=Sage}
\setbeamercolor{date}{fg=StoneLight}
\setbeamercolor{frametitle}{fg=Slate, bg=}
\setbeamercolor{framesubtitle}{fg=Sage}
\setbeamercolor{itemize item}{fg=Slate}
\setbeamercolor{itemize subitem}{fg=Sage}
\setbeamercolor{enumerate item}{fg=Slate}
\setbeamercolor{section in toc}{fg=Stone}
\setbeamercolor{bibliography entry author}{fg=Stone}
\setbeamercolor{bibliography entry title}{fg=Stone}
\setbeamercolor{bibliography entry note}{fg=StoneLight}

\setbeamerfont{title}{size=\Large, series=\bfseries}
\setbeamerfont{subtitle}{size=\normalsize}
\setbeamerfont{frametitle}{size=\large, series=\bfseries}
\setbeamerfont{framesubtitle}{size=\small}
\setbeamerfont{author}{size=\normalsize}
\setbeamerfont{institute}{size=\small}
\setbeamerfont{date}{size=\small}
\setbeamerfont{footline}{size=\tiny}

% Puces : petits carrés arrondis ; numéros : badges
\setbeamertemplate{itemize item}{\tikz[baseline=-0.6ex]\fill[Slate, rounded corners=1pt] (0,0) rectangle (0.32em,0.32em);}
\setbeamertemplate{itemize subitem}{\tikz[baseline=-0.5ex]\draw[Sage, thick, rounded corners=0.5pt] (0,0) rectangle (0.26em,0.26em);}
\setbeamertemplate{enumerate items}{\tikz[baseline=-0.7ex]{\node[fill=Slate, text=SoftWhite, rounded corners=2.5pt,
  inner sep=2pt, minimum width=1.3em, minimum height=1.3em, font=\footnotesize\bfseries] {\insertenumlabel};}}

% --- Page de titre ---
\defbeamertemplate*{title page}{medical}{
  \vspace{1.2cm}
  \hspace{0.8cm}%
  \begin{minipage}[t]{0.78\textwidth}
    {\usebeamerfont{title}\usebeamercolor[fg]{title}\inserttitle\par}
    \vspace{0.25cm}
    {\lightfont\usebeamercolor[fg]{subtitle}\insertsubtitle\par}
    \vspace{0.8cm}
    {\color{WarmAmber}\rule{2.2cm}{1.5pt}\par}
    \vspace{0.6cm}
    {\mediumfont\usebeamercolor[fg]{author}\insertauthor\par}
    \vspace{0.1cm}
    {\lightfont\small\usebeamercolor[fg]{institute}\insertinstitute\par}
    \vspace{0.25cm}
    {\lightfont\small\usebeamercolor[fg]{date}\insertdate\par}
  \end{minipage}}

% --- Titre de diapositive ---
\makeatletter
\defbeamertemplate*{frametitle}{medical}{
  \vspace{0.02cm}
  \hspace{0.15cm}%
  \begin{minipage}[t]{0.92\textwidth}
    \usebeamerfont{frametitle}\usebeamercolor[fg]{frametitle}\insertframetitle%
    \ifx\insertframesubtitle\@empty\else
      \\[1pt]{\lightfont\small\usebeamercolor[fg]{framesubtitle}\insertframesubtitle}
    \fi
  \end{minipage}
  \vspace{-0.05cm}}
\makeatother

% --- Pied de page (réserve sa hauteur, puis se dessine en surimpression) ---
\defbeamertemplate*{footline}{medical}{
  \hbox to \paperwidth{\rule{0pt}{0.6cm}\hfill}%
  \begin{tikzpicture}[remember picture, overlay]
    \fill[Pearl] (current page.south west) rectangle ([yshift=0.45cm]current page.south east);
    \node[anchor=west, inner sep=0pt] at ([xshift=0.5cm, yshift=0.22cm]current page.south west)
      {\usebeamerfont{footline}\color{StoneLight}\insertshortauthor{} \textbar{} \insertshortinstitute};
    \node[anchor=center, inner sep=0pt] at ([yshift=0.22cm]current page.south)
      {\usebeamerfont{footline}\color{StoneLight}\insertshorttitle};
    \node[anchor=east, inner sep=0pt] at ([xshift=-0.5cm, yshift=0.22cm]current page.south east)
      {\usebeamerfont{footline}\color{Sage}\bfseries\insertframenumber\,/\,\inserttotalframenumber};
  \end{tikzpicture}}

% --- Fond des pages spéciales : barre latérale + forme optionnelle ---
\newcommand{\fondbarre}[1]{%
  \setbeamertemplate{background canvas}{%
    \begin{tikzpicture}[x=\paperwidth, y=\paperheight]
      \useasboundingbox (0,0) rectangle (1,1);
      \fill[SoftWhite] (0,0) rectangle (1,1);
      \fill[Slate] (0,0) rectangle (0.043,1);
      #1
    \end{tikzpicture}}}

% Projection radiographique d'une sphère (page de titre)
\newcommand{\formesphere}{%
  \coordinate (sC) at (0.86,0.77);
  \foreach \sc/\op in {1.0/0.06, 0.82/0.06, 0.64/0.07, 0.46/0.08, 0.28/0.1, 0.12/0.12} {
    \pgfmathsetmacro{\rr}{3.2*\sc}
    \fill[Slate, opacity=\op] (sC) circle (\rr cm);}}

% Projection radiographique d'un cube vu selon sa diagonale (page de fin)
\newcommand{\formecube}{%
  \coordinate (hC) at (0.89,0.72);
  \pgfmathsetmacro{\cubeS}{2.2}\pgfmathsetmacro{\cA}{cos(30)}\pgfmathsetmacro{\sA}{sin(30)}%
  \pgfmathsetmacro{\cB}{cos(45)}\pgfmathsetmacro{\sB}{sin(45)}%
  \pgfmathsetmacro{\aV}{\cubeS*(\cA+\sA)}\pgfmathsetmacro{\bV}{\cubeS*(\cA-\sA)}%
  \pgfmathsetmacro{\uV}{\bV*\cB}\pgfmathsetmacro{\vV}{\aV*\cB}\pgfmathsetmacro{\hV}{\cubeS*\sB}%
  \foreach \sc/\op in {1.0/0.06, 0.84/0.06, 0.68/0.07, 0.52/0.08, 0.36/0.1, 0.2/0.12} {
    \fill[Slate, opacity=\op]
      ([shift={(\sc*\aV cm, \sc*(-\uV+\hV) cm)}]hC) -- ([shift={(\sc*\bV cm, \sc*(\vV+\hV) cm)}]hC) --
      ([shift={(-\sc*\aV cm, \sc*(\uV+\hV) cm)}]hC) -- ([shift={(-\sc*\aV cm, \sc*(\uV-\hV) cm)}]hC) --
      ([shift={(-\sc*\bV cm, \sc*(-\vV-\hV) cm)}]hC) -- ([shift={(\sc*\aV cm, \sc*(-\uV-\hV) cm)}]hC) -- cycle;}}

% \pagetitre et \pagefin : à appeler dans le document
\newcommand{\pagetitre}{{%
  \setbeamertemplate{footline}{}\fondbarre{\formesphere}
  \begin{frame}[plain]\titlepage\end{frame}}}
\newcommand{\pagefin}[1]{{%   #1 : ligne de contact
  \setbeamertemplate{footline}{}\fondbarre{\formecube}
  \begin{frame}[plain]
    \vspace{2.5cm}\hspace{0.8cm}%
    \begin{minipage}{0.78\textwidth}
      {\color{Slate}\LARGE\bfseries Merci de votre attention}\\[0.35cm]
      {\lightfont\color{Sage}\normalsize #1}\\[0.15cm]
      \vspace{0.6cm}{\color{WarmAmber}\rule{2.2cm}{1.5pt}}
    \end{minipage}
  \end{frame}}}

% --- Page de section automatique, avec grand numéro ---
\AtBeginSection[]{{%
  \setbeamertemplate{footline}{}
  \fondbarre{\node[anchor=north east, font=\fontsize{100}{100}\selectfont\bfseries, text=SlatePale] at (0.93,0.95) {\thesection};}
  \begin{frame}[plain]
    \hypertarget{section\thesection}{}
    \vspace{2.5cm}\hspace{1.1cm}%
    \begin{minipage}{0.75\textwidth}\raggedright
      {\lightfont\color{Sage}\large Partie \thesection}\\[0.35cm]
      {\color{Slate}\LARGE\bfseries\insertsectionhead}\\[0.3cm]
      {\color{WarmAmber}\rule{1.8cm}{1.5pt}}
    \end{minipage}
  \end{frame}}}

% --- Sommaire : entrées cliquables à badges ---
\setbeamertemplate{section in toc}{%
  \leavevmode
  \hyperlink{section\inserttocsectionnumber}{%
    \tikz[baseline=-0.7ex]{\node[fill=Slate, text=SoftWhite, rounded corners=3pt, inner sep=3pt,
      minimum width=1.5em, minimum height=1.5em, font=\small\bfseries] {\inserttocsectionnumber};}%
    \hspace{0.4em}{\mediumfont\color{Stone}\inserttocsection}}\par
  \vspace{0.3em}}

% --- Aides de mise en page ---
\newcommand{\twocols}[2]{\begin{columns}[T]\begin{column}{0.47\textwidth}#1\end{column}\begin{column}{0.47\textwidth}#2\end{column}\end{columns}}
\newcommand{\highlight}[1]{\tikz[baseline=-0.7ex]{\node[fill=SlatePale, text=SlateDeep, rounded corners=3pt, inner xsep=5pt, inner ysep=2pt, font=\small] {#1};}}
% Source d'une image, en petit sous la figure
\newcommand{\sourceimage}[1]{\par{\tiny\color{StoneLight}#1\par}}
